% ═══════════════════════════════════════════════════════════════
% MUSTERLÖSUNG: Redoxreaktionen komplett
% Thema 2.5 Redoxreaktionen | Chemie Klasse 8 | 42 BE
% ═══════════════════════════════════════════════════════════════

\documentclass[10pt, a4paper]{article}

\usepackage[utf8]{inputenc}
\usepackage[T1]{fontenc}
\usepackage[ngerman]{babel}

\usepackage{helvet}
\renewcommand{\familydefault}{\sfdefault}

\usepackage{microtype}
\usepackage[margin=1.2cm, top=1.6cm, bottom=1.2cm]{geometry}
\usepackage{enumitem}
\usepackage{tabularx}
\usepackage{booktabs}
\usepackage{array}
\usepackage{multicol}
\usepackage{tikz}
\usetikzlibrary{calc, shapes.geometric, arrows.meta, positioning}

\usepackage{amsmath, amssymb}
\usepackage{siunitx}
\sisetup{locale = DE, per-mode = symbol}

\usepackage{fancyhdr}
\usepackage{lastpage}

\usepackage{xcolor}
\usepackage{tcolorbox}
\tcbuselibrary{skins}

\usepackage{qrcode}
\usepackage[hidelinks]{hyperref}

% ═══════════════════════════════════════════════════════════════
% FARBEN
% ═══════════════════════════════════════════════════════════════

\definecolor{chemie}{HTML}{2a7a4b}
\definecolor{hellgrau}{HTML}{F5F5F5}
\definecolor{mittelgrau}{HTML}{888888}
\definecolor{dunkelgrau}{HTML}{444444}
\definecolor{silber}{HTML}{C0C0C0}
\definecolor{kupfer}{HTML}{B87333}
\definecolor{sauerstoff}{HTML}{c62828}
\definecolor{warnung}{HTML}{b35c00}
\definecolor{loesung}{HTML}{c62828}

\colorlet{fachfarbe}{chemie}

% ═══════════════════════════════════════════════════════════════
% VARIABLEN
% ═══════════════════════════════════════════════════════════════

\newcommand{\fach}{Chemie}
\newcommand{\thema}{Redoxreaktionen komplett}
\newcommand{\klasse}{8}
\newcommand{\maxpunkte}{42}

% ═══════════════════════════════════════════════════════════════
% KOPF- UND FUSSZEILE
% ═══════════════════════════════════════════════════════════════

\pagestyle{fancy}
\fancyhf{}
\renewcommand{\headrulewidth}{0.4pt}
\renewcommand{\footrulewidth}{0pt}
\lhead{\small \fach{} Klasse \klasse{} | \thema{} | \textcolor{loesung}{\textbf{MUSTERLÖSUNG}}}
\rhead{}
\cfoot{\small\textcolor{mittelgrau}{Seite \thepage{} von \pageref{LastPage}}}

% ═══════════════════════════════════════════════════════════════
% BOXEN
% ═══════════════════════════════════════════════════════════════

\newtcolorbox{merkkasten}[1][]{
    colback=hellgrau,
    colframe=hellgrau,
    boxrule=0pt,
    arc=0pt,
    left=5pt, right=5pt, top=4pt, bottom=3pt,
    borderline north={2pt}{0pt}{fachfarbe},
    before skip=3pt, after skip=3pt,
    title={#1},
    fonttitle=\bfseries\small,
    coltitle=fachfarbe,
    attach boxed title to top left={yshift=-2mm, xshift=4mm},
    boxed title style={colback=hellgrau, arc=0pt, boxrule=0pt}
}

\newtcolorbox{sicherheitsbox}{
    colback=warnung!10,
    colframe=warnung,
    boxrule=1.5pt,
    arc=0pt,
    left=6pt, right=6pt, top=6pt, bottom=4pt,
    before skip=4pt, after skip=4pt
}

\newtcolorbox{teilbox}[1][]{
    colback=fachfarbe!8,
    colframe=fachfarbe,
    boxrule=0pt,
    borderline west={3pt}{0pt}{fachfarbe},
    arc=0pt,
    left=6pt, right=6pt, top=4pt, bottom=3pt,
    before skip=6pt, after skip=3pt
}

% ═══════════════════════════════════════════════════════════════
% BEFEHLE
% ═══════════════════════════════════════════════════════════════

\newcommand{\lsg}[1]{\textcolor{loesung}{\textbf{#1}}}
\newcommand{\punktebox}[1]{\hfill\small\textcolor{mittelgrau}{[#1/#1]}}

\newcommand{\kasten}[2]{%
    \vspace{4pt}%
    \noindent\textbf{\textcolor{fachfarbe}{Kasten #1:}} #2%
    \vspace{2pt}%
    \nopagebreak[4]%
}

\newcommand{\operator}[1]{\textbf{\textcolor{fachfarbe}{#1}}}

\setlength{\parindent}{0pt}
\setlength{\parskip}{1pt}

% ═══════════════════════════════════════════════════════════════
% DOKUMENT
% ═══════════════════════════════════════════════════════════════

\begin{document}

% ═══════════════════════════════════════════════════════════════
% HEADER
% ═══════════════════════════════════════════════════════════════

\begin{center}
    {\large\bfseries\textcolor{fachfarbe}{\thema}}\\[0.1em]
    {\small \textcolor{loesung}{\textbf{MUSTERLÖSUNG}} | \fach{} Klasse \klasse{} | Thema 2.5 Redoxreaktionen}
\end{center}
\vspace{-0.4em}
\hrule
\vspace{0.3em}

% ═══════════════════════════════════════════════════════════════
% TEIL 1: REAKTIONSGLEICHUNGEN
% ═══════════════════════════════════════════════════════════════

\begin{teilbox}
\textbf{\textcolor{fachfarbe}{Teil 1: Reaktionsgleichungen}}
\end{teilbox}

\vspace{-2pt}

% KASTEN 1: Begriffe
\begin{merkkasten}[Kasten 1: Begriffe einer Reaktionsgleichung]
\punktebox{3}

\vspace{3pt}
\begin{tabular}{@{}p{3cm}p{4cm}p{6.5cm}@{}}
\textbf{Begriff} & \textbf{Bedeutung} & \textbf{Position in der Gleichung} \\[0.3em]
\lsg{Edukte} & Ausgangsstoffe & \lsg{links} vom Pfeil \\[0.6em]
\lsg{Produkte} & Endstoffe & \lsg{rechts} vom Pfeil \\[0.6em]
$\longrightarrow$ & \lsg{„reagiert zu"} & in der Mitte
\end{tabular}
\end{merkkasten}

% KASTEN 2: Wortgleichung
\begin{merkkasten}[Kasten 2: Wortgleichung]
\punktebox{2}

\vspace{4pt}
\hspace{1.5em} \lsg{Magnesium} + \lsg{Sauerstoff} $\longrightarrow$ \lsg{Magnesiumoxid}
\end{merkkasten}

% KASTEN 3: Formelgleichungen
\begin{merkkasten}[Kasten 3: Formelgleichungen]
\punktebox{6}

\vspace{4pt}
\textbf{Magnesium:} \hspace{1em} \lsg{2} Mg + O$_2$ $\longrightarrow$ \lsg{2} MgO

\vspace{4pt}
\textbf{Calcium:} \hspace{1.5em} \lsg{2} Ca + O$_2$ $\longrightarrow$ \lsg{2} CaO

\vspace{4pt}
\textbf{Natrium:} \hspace{1.5em} \lsg{4} Na + O$_2$ $\longrightarrow$ \lsg{2} Na$_2$O
\end{merkkasten}

\vspace{2pt}

% ═══════════════════════════════════════════════════════════════
% TEIL 2: REDUKTION
% ═══════════════════════════════════════════════════════════════

\begin{teilbox}
\textbf{\textcolor{fachfarbe}{Teil 2: Reduktion – Silberoxid zerlegen}}
\end{teilbox}

\vspace{-2pt}

% KASTEN 4: Definition Reduktion
\begin{merkkasten}[Kasten 4: Definition Reduktion]
\textbf{Reduktion} kommt vom lat. \textit{reducere} = \lsg{re} + \lsg{ducere} = \lsg{zurückführen} \punktebox{2}

\vspace{3pt}
\textbf{Definition:} Reduktion ist der \lsg{Entzug von Sauerstoff} aus einer Verbindung.
\end{merkkasten}

% KASTEN 5: Sicherheit
\begin{sicherheitsbox}
\textbf{\textcolor{warnung}{Kasten 5: Sicherheitsregeln}} \punktebox{4}

\vspace{2pt}
\begin{minipage}[t]{0.6\textwidth}
\begin{itemize}[leftmargin=1.3em, itemsep=3pt]
    \item[\textcolor{warnung}{\textbf{!}}] \lsg{Schutzbrille} tragen
    \item[\textcolor{warnung}{\textbf{!}}] Im \lsg{Abzug} arbeiten
    \item[\textcolor{warnung}{\textbf{!}}] Heißes Glas \lsg{nicht anfassen}
    \item[\textcolor{warnung}{\textbf{!}}] \lsg{Reagenzglashalter} verwenden
\end{itemize}
\end{minipage}%
\hfill
\begin{minipage}[t]{0.35\textwidth}
\centering
\textbf{Material:}\\
{\footnotesize Silberoxid (Ag$_2$O), Reagenzglas, Brenner, Glimmspan}
\end{minipage}
\end{sicherheitsbox}

% KASTEN 6: Beobachtung
\kasten{6}{Beobachtung – Silberoxid erhitzen} \punktebox{3}

\vspace{2pt}
Das \textbf{schwarze} Silberoxid wird beim Erhitzen \lsg{silbrig-glänzend}.

Der glühende Holzspan \lsg{flammt auf}. $\Rightarrow$ Nachweis für \lsg{O$_2$} (Glimmspanprobe).

\vspace{4pt}

% KASTEN 7: Reaktionsgleichung Ag2O
\begin{merkkasten}[Kasten 7: Reaktionsgleichung Silberoxid]
\textbf{Wortgleichung:} \lsg{Silberoxid} $\longrightarrow$ \lsg{Silber} + \lsg{Sauerstoff} \punktebox{3}

\vspace{3pt}
\textbf{Formelgleichung:} \lsg{2} Ag$_2$O $\longrightarrow$ \lsg{4} Ag + \lsg{1} O$_2$
\end{merkkasten}

% KASTEN 8: Vergleich
\kasten{8}{Vergleich: Oxidation und Reduktion} \punktebox{4}

\vspace{2pt}
\begin{tabular}{|p{3.5cm}|p{5cm}|p{5cm}|}
\hline
& \textbf{Oxidation} & \textbf{Reduktion} \\
\hline
\textbf{Sauerstoff wird...} & \lsg{aufgenommen} & \lsg{abgegeben} \\
\hline
\textbf{Merkhilfe} & \lsg{„O kommt dazu"} & \lsg{„O geht weg"} \\
\hline
\textbf{Beispiel (Silber)} & 4 Ag + O$_2$ $\rightarrow$ 2 Ag$_2$O & 2 Ag$_2$O $\rightarrow$ 4 Ag + O$_2$ \\
\hline
\end{tabular}

\vspace{4pt}

% ═══════════════════════════════════════════════════════════════
% TEIL 3: REDOXREAKTION
% ═══════════════════════════════════════════════════════════════

\begin{teilbox}
\textbf{\textcolor{fachfarbe}{Teil 3: Redoxreaktion – Sauerstoffübertragung}}
\end{teilbox}

\vspace{-2pt}

% KASTEN 9: Beobachtung CuO
\kasten{9}{Beobachtung – CuO + C erhitzen} \punktebox{3}

\vspace{2pt}
Das \textbf{schwarze} Gemisch (Kupferoxid + Kohlenstoff) wird beim Erhitzen \lsg{rötlich (kupferfarben)}.

Es entsteht ein farbloses Gas und zurück bleibt: \lsg{reines Kupfer (Cu)}

\vspace{4pt}

% KASTEN 10: Redox-Schema
\begin{merkkasten}[Kasten 10: Redox-Schema – Sauerstoffübertragung]
\punktebox{4}

\vspace{6pt}
\begin{center}
\begin{tikzpicture}[scale=1]
    % Edukte
    \node[draw, fill=hellgrau, minimum width=2cm, minimum height=0.8cm] (cuo) at (0,0) {CuO};
    \node[below=0.1cm of cuo, font=\scriptsize] {Kupferoxid};

    \node at (1.5,0) {$+$};

    \node[draw, fill=hellgrau, minimum width=1.2cm, minimum height=0.8cm] (c) at (2.8,0) {C};
    \node[below=0.1cm of c, font=\scriptsize] {Kohlenstoff};

    \node at (4.3,0) {$\longrightarrow$};

    % Produkte
    \node[draw, fill=kupfer!30, minimum width=1.2cm, minimum height=0.8cm] (cu) at (5.8,0) {Cu};
    \node[below=0.1cm of cu, font=\scriptsize] {Kupfer};

    \node at (7,0) {$+$};

    \node[draw, fill=hellgrau, minimum width=1.5cm, minimum height=0.8cm] (co2) at (8.5,0) {CO$_2$};
    \node[below=0.1cm of co2, font=\scriptsize] {Kohlenstoffdioxid};

    % O-Übertragungspfeil (LÖSUNG)
    \draw[->, very thick, loesung] (0,0.7) .. controls (2,1.5) and (6.5,1.5) .. (8.5,0.7);
    \node[above, font=\scriptsize\bfseries, loesung] at (4.25,1.4) {O};
\end{tikzpicture}
\end{center}

\vspace{3pt}
CuO wird \lsg{reduziert} (gibt O ab). \hfill C wird \lsg{oxidiert} (nimmt O auf).
\end{merkkasten}

% KASTEN 11: Ox-/Red-Mittel + Thermit
\kasten{11}{Oxidations- und Reduktionsmittel + Thermit} \punktebox{4}

\vspace{2pt}
\begin{minipage}[t]{0.5\textwidth}
\begin{tabular}{|p{3.2cm}|p{3cm}|}
\hline
\textbf{Begriff} & \textbf{Beispiel (CuO+C)} \\
\hline
\textbf{Oxidationsmittel}\newline{\scriptsize(gibt O ab)} & \lsg{CuO} \\
\hline
\textbf{Reduktionsmittel}\newline{\scriptsize(nimmt O auf)} & \lsg{C} \\
\hline
\end{tabular}
\end{minipage}%
\hfill
\begin{minipage}[t]{0.45\textwidth}
\textbf{Thermit-Reaktion:}

\vspace{3pt}
Fe$_2$O$_3$ + \lsg{2} Al $\longrightarrow$ \lsg{2} Fe + Al$_2$O$_3$

\vspace{3pt}
{\small Anwendung: \lsg{Schweißen von Schienen}}
\end{minipage}

\vspace{4pt}

% ═══════════════════════════════════════════════════════════════
% TEIL 4: SICHERUNG
% ═══════════════════════════════════════════════════════════════

\begin{teilbox}
\textbf{\textcolor{fachfarbe}{Teil 4: Zusammenfassung}}
\end{teilbox}

\vspace{-2pt}

% KASTEN 12: Definition + Bedeutung
\begin{merkkasten}[Kasten 12: Definition Redoxreaktion und Bedeutung]
\punktebox{4}

\vspace{3pt}
\textbf{\textcolor{fachfarbe}{Redoxreaktion}} = Reaktion mit \lsg{Sauerstoffübertragung}

\vspace{3pt}
Dabei laufen \lsg{Oxidation} und \lsg{Reduktion} \textbf{gleichzeitig} ab.

\vspace{6pt}
\textbf{Bedeutung:} Redoxreaktionen sind die Grundlage der \lsg{Metallgewinnung}!

\vspace{3pt}
Beispiele: \lsg{Hochofen (Eisen), Thermit (Schweißen), Kupfergewinnung}
\end{merkkasten}

% ═══════════════════════════════════════════════════════════════
% GESAMTPUNKTE
% ═══════════════════════════════════════════════════════════════

\vfill

\begin{center}
{\large\textbf{Gesamt: \maxpunkte{} / \maxpunkte{} BE}}
\end{center}

\end{document}
